% Part opener. Structural figures get their own number prefix so they
% do not inherit a misleading chapter number.
\structuralfigures{P4.}

\vspace*{0.4cm}
Parts~I to~III built a network that forwards packets. Nobody uses a network like
that. Users expect an address to appear without asking for one, a name to resolve
to that address, a private network to reach the internet through one public
address, and somebody to have decided which of them is allowed where.

That is this part. It is the largest single block of the book because it carries
two blueprint domains at once: all seven topics of domain~4, plus topic~\tp{1.7}
from domain~1, which sits here rather than in Part~I because troubleshooting DHCP
is impossible before you know what a relay is doing.

\ccnafigh{%
\begin{tikzpicture}[font=\scriptsize]
  \tikzset{
    stg/.style={rounded corners=3pt, draw=#1, fill=#1!8, text=#1!45!black,
                line width=1pt, text width=3.1cm, align=center,
                minimum height=1.5cm, font=\scriptsize\bfseries},
    band/.style={font=\scriptsize\bfseries, text=#1!45!black, inner sep=1pt},
    fl/.style={-{Stealth[length=2.4mm]}, line width=0.9pt, neutral!70}}

  % Band 1 -- services. Boxes start at x=2.2 to leave the rotated band
  % label clear of the first box's left edge, which sits 1.7cm out.
  \node[stg=dom1] (dhcp) at (2.2,0)  {\faIcon{address-card}\\[2pt]20. DHCPv4\\
      \footnotesize\mdseries an address appears};
  \node[stg=dom4] (dns)  at (6.2,0)  {\faIcon{book}\\[2pt]21. DNS records\\
      \footnotesize\mdseries a name resolves};
  \node[stg=dom4] (nat)  at (10.2,0) {\faIcon{exchange-alt}\\[2pt]22. NAT and PAT\\
      \footnotesize\mdseries to the internet};

  % Band 2 -- security
  \node[stg=dom4] (aaa)  at (2.2,-2.6)  {\faIcon{user-lock}\\[2pt]23. AAA, SFTP\\
      \footnotesize\mdseries who may log in};
  \node[stg=dom4] (acl)  at (6.2,-2.6)  {\faIcon{filter}\\[2pt]24. ACLs\\
      \footnotesize\mdseries what may pass};
  \node[stg=dom4] (l2s)  at (10.2,-2.6) {\faIcon{shield-alt}\\[2pt]25. Layer 2 security\\
      \footnotesize\mdseries the access port};
  \node[stg=dom4] (vpn)  at (14.2,-2.6) {\faIcon{lock}\\[2pt]26. IPsec VPNs\\
      \footnotesize\mdseries across the internet};

  \draw[fl] (dhcp) -- (dns);
  \draw[fl] (dns)  -- (nat);
  \draw[fl] (aaa)  -- (acl);
  \draw[fl] (acl)  -- (l2s);
  \draw[fl] (l2s)  -- (vpn);

  \node[band=dom1, rotate=90] at (0.2,0)    {SERVICES};
  \node[band=dom4, rotate=90] at (0.2,-2.6) {SECURITY};
\end{tikzpicture}}%
{Part~IV. Three chapters make the network usable; four decide who is allowed to
use it. Chapter~20 belongs to domain~1 and the rest to domain~4.}{p4map}

\begin{examnote}[Six verbs, and only one of them is \emph{describe}]
Topic~\tp{1.7} says \emph{troubleshoot}. Topics \tp{4.1}, \tp{4.3}, \tp{4.6} and
\tp{4.7} say \emph{configure}. Topic~\tp{4.2} says \emph{manage} and \tp{4.4} says
\emph{diagnose}. Only \tp{4.5}, the VPN topic, says \emph{describe} --- which is
why Chapter~26 is the shortest chapter in the part and has no lab.

Read that distribution as an instruction. Four of these seven chapters will be
examined by putting a configuration in front of you and asking what it does or
what is wrong with it. Reading an ACL correctly is worth more marks here than
being able to recite what a wildcard mask is.
\end{examnote}

\begin{conceptbox}[Two platforms from here on]
Parts~I--III were written for Packet Tracer, which models switching and routing
well enough. From this part onward the labs move to \term{CML} or \term{GNS3}
where the platform matters: topic~\tp{4.3} names \textbf{IOS XE} explicitly, and
SFTP, SCP and AAA against a real RADIUS server need a control plane that Packet
Tracer only pretends to have.

Every lab header names its platform. Where a lab can still be done usefully in
Packet Tracer with a caveat, the caveat is stated rather than hidden.
\end{conceptbox}
